% unit: u03
% sources: sources/072-elimination-with-matrices.txt, sources/123-lecture-elimination-with-matrices.txt, sources/124-mit18-06scf11-ses1-2sum-pdf-download-file.txt, sources/026-elimination-with-matrices.txt, sources/125-mit18-06scf11-ses1-2prob-pdf-download-file.txt, sources/126-mit18-06scf11-ses1-2sol-pdf-download-file.txt
\section{Elimination with matrices}\label{sec:u03}

Section~\ref{sec:u01} asked which combination of the columns produces $\vb{b}$ but had no
systematic way to find it. \emph{Gaussian elimination} is that way, and it is how essentially
every software package solves $A\vb{x}=\vb{b}$. It subtracts multiples of equations from
other equations until the system is triangular, then solves the triangular system from the
bottom up. This section does elimination by hand, sees how it can fail, and then rewrites
each step as multiplication by a simple matrix. That last move turns a procedure into
algebra, and it is the starting point for the factorization $A=LU$ in Section~\ref{sec:u05}.

\subsection{Pivots and multipliers}

Our example has three equations in three unknowns:
\begin{equation}\label{eq:u03-system}
\begin{aligned}
  x + 2y + z &= 2\\
  3x + 8y + z &= 12\\
  4y + z &= 2
\end{aligned}
\qquad\qquad
A=\begin{bmatrix} 1 & 2 & 1\\ 3 & 8 & 1\\ 0 & 4 & 1\end{bmatrix},\quad
\vb{b}=\begin{bmatrix} 2\\ 12\\ 2\end{bmatrix}.
\end{equation}
Elimination works on the rows of $A$ (the equations), so we write only the numbers.

\begin{definition}[Pivot, multiplier, upper triangular]\label{def:u03-pivot}
Elimination clears the columns of $A$ from left to right. At step $k$ the \emph{pivot} is the
entry in position $(k,k)$ of the current matrix; it must be nonzero. To clear an entry
$a_{ik}$ below the pivot ($i>k$), subtract
\[
  \ell_{ik} = \frac{\text{entry to eliminate}}{\text{pivot}}
\]
times the pivot row $k$ from row $i$. The number $\ell_{ik}$ is the \emph{multiplier}. When
all columns are cleared the result $U$ is \emph{upper triangular}: every entry below the main
diagonal is zero, and the pivots sit on the diagonal.
\end{definition}

\begin{example}[Forward elimination on a $3\times3$ system]\label{ex:u03-forward}
Eliminate below the pivots of the matrix $A$ in \eqref{eq:u03-system}.
\begin{itemize}
  \item \emph{Step $(2,1)$.} The first pivot is $1$. The entry below it is $3$, so
    $\ell_{21}=3/1=3$: subtract $3$ times row $1$ from row $2$, giving
    $(3,8,1)-3(1,2,1)=(0,2,-2)$.
  \item \emph{Step $(3,1)$.} The entry in position $(3,1)$ is already $0$, so
    $\ell_{31}=0$ and nothing changes. (Software still checks this position.)
  \item \emph{Step $(3,2)$.} The second pivot is the $2$ now in position $(2,2)$. The entry
    below it is $4$, so $\ell_{32}=4/2=2$: subtract $2$ times the \emph{new} row $2$ from
    row $3$, giving $(0,4,1)-2(0,2,-2)=(0,0,5)$.
\end{itemize}
The third pivot is $5$, and
\[
  A=\begin{bmatrix} 1 & 2 & 1\\ 3 & 8 & 1\\ 0 & 4 & 1\end{bmatrix}
  \;\longrightarrow\;
  \begin{bmatrix} 1 & 2 & 1\\ 0 & 2 & -2\\ 0 & 4 & 1\end{bmatrix}
  \;\longrightarrow\;
  U=\begin{bmatrix} 1 & 2 & 1\\ 0 & 2 & -2\\ 0 & 0 & 5\end{bmatrix}.
\]
\end{example}
The pivots $1,2,5$ already tell us a great deal about $A$. For instance, their product
$1\cdot2\cdot5=10$ is the determinant of $A$ (Section~\ref{sec:u20}).

The right-hand side must receive the same operations. The convenient bookkeeping is the
\emph{augmented matrix} $[\,A\;\;\vb{b}\,]$, with $\vb{b}$ tacked on as an extra column.
Subtracting $3$ times row~1 changes $12$ into $12-3\cdot2=6$; subtracting $2$ times row~2
changes $2$ into $2-2\cdot6=-10$. Elimination turns $A\vb{x}=\vb{b}$ into $U\vb{x}=\vb{c}$
with $\vb{c}=(2,6,-10)$; Figure~\ref{fig:u03-flow} records the whole computation.

\begin{figure}[ht]
\centering
\small
\begin{tikzpicture}[>=Stealth]
  \node (A) {$\left[\begin{array}{rrr|r} \boxed{1} & 2 & 1 & 2\\ 3 & 8 & 1 & 12\\ 0 & 4 & 1 & 2\end{array}\right]$};
  \node[right=2.3cm of A] (B) {$\left[\begin{array}{rrr|r} \boxed{1} & 2 & 1 & 2\\ 0 & \boxed{2} & -2 & 6\\ 0 & 4 & 1 & 2\end{array}\right]$};
  \node[right=2.3cm of B] (C) {$\left[\begin{array}{rrr|r} \boxed{1} & 2 & 1 & 2\\ 0 & \boxed{2} & -2 & 6\\ 0 & 0 & \boxed{5} & -10\end{array}\right]$};
  \draw[->,thick] (A) -- node[above,align=center,font=\footnotesize] {$E_{21}$\\ row$_2-3\,$row$_1$} (B);
  \draw[->,thick] (B) -- node[above,align=center,font=\footnotesize] {$E_{32}$\\ row$_3-2\,$row$_2$} (C);
  \node[below=3pt of A] {$[\,A\;\;\vb{b}\,]$};
  \node[below=3pt of B] {$[\,E_{21}A\;\;E_{21}\vb{b}\,]$};
  \node[below=3pt of C] {$[\,U\;\;\vb{c}\,]$};
\end{tikzpicture}
\caption{Forward elimination on the augmented matrix. Boxed entries are pivots. Each arrow is
one row operation, equivalently one multiplication on the left by an elimination matrix; the
step $E_{31}$ is skipped because the $(3,1)$ entry is already zero. The right-hand column
changes from $\vb{b}=(2,12,2)$ to $\vb{c}=(2,6,-10)$.}\label{fig:u03-flow}
\end{figure}

\subsection{Back substitution}

The triangular system $U\vb{x}=\vb{c}$ reads
\[
  x+2y+z=2,\qquad 2y-2z=6,\qquad 5z=-10 .
\]
Solve it in reverse order. The last equation gives $z=-2$. The second becomes
$2y+4=6$, so $y=1$. The first becomes $x+2-2=2$, so $x=2$. Checking in the original
equations \eqref{eq:u03-system}: $2+2-2=2$, $6+8-2=12$, $4-2=2$.

\begin{proposition}[Back substitution]\label{prop:u03-back}
Let $U$ be an $n\times n$ upper triangular matrix whose diagonal entries $u_{11},\dots,u_{nn}$
are all nonzero. Then for every $\vb{c}$ the system $U\vb{x}=\vb{c}$ has exactly one
solution, given by
\[
  x_n=\frac{c_n}{u_{nn}},\qquad
  x_i=\frac{1}{u_{ii}}\Bigl(c_i-\sum_{j=i+1}^{n}u_{ij}x_j\Bigr)\quad (i=n-1,\dots,1).
\]
\end{proposition}

\begin{proof}
Equation $i$ of $U\vb{x}=\vb{c}$ is $u_{ii}x_i+\sum_{j>i}u_{ij}x_j=c_i$, because the entries
left of the diagonal are zero. Equation $n$ involves only $x_n$ and forces $x_n=c_n/u_{nn}$.
Suppose $x_{i+1},\dots,x_n$ are already forced. Then equation $i$ has exactly one unknown
left, $x_i$, with the nonzero coefficient $u_{ii}$, so it forces the stated value. Working
upward, every component is forced, which proves uniqueness; and the values so obtained
satisfy every equation, which proves existence.
\end{proof}

\subsection{When elimination breaks down}

Zero is never allowed as a pivot. If a zero appears in the pivot position, there are two
cases.

\emph{Temporary failure.} If some row \emph{below} has a nonzero entry in that column,
exchange the two rows and continue.

\begin{example}[A row exchange rescues elimination]\label{ex:u03-exchange}
Change the $8$ in $A$ to a $6$. The first step gives
row~2 $=(3,6,1)-3(1,2,1)=(0,0,-2)$, a zero in the $(2,2)$ pivot position. Row~3 is
$(0,4,1)$, with a nonzero entry there, so exchange rows~2 and~3:
\[
  \begin{bmatrix} 1 & 2 & 1\\ 0 & 0 & -2\\ 0 & 4 & 1\end{bmatrix}
  \;\xrightarrow{\text{exchange rows 2, 3}}\;
  \begin{bmatrix} 1 & 2 & 1\\ 0 & 4 & 1\\ 0 & 0 & -2\end{bmatrix},
\]
which is already triangular, with pivots $1,4,-2$. The matrix is fine; only the order of the
equations was unlucky.
\end{example}

\emph{Complete failure.} If the pivot position holds a zero and every entry below it is zero
too, no exchange helps.

\begin{example}[Elimination fails: a singular matrix]\label{ex:u03-singular}
Change the $(3,3)$ entry of $A$ from $1$ to $-4$. The first step
gives row~2 $=(0,2,-2)$ and leaves row~3 $=(0,4,-4)$; subtracting $2$ times row~2 leaves
row~3 $=(0,0,0)$. There is no third pivot. The matrix is \emph{singular}: depending on
$\vb{b}$, the system has no solution or infinitely many, never exactly one. (In the column
picture, column~3 has become $(1,1,-4)=3\,(1,3,0)-(2,8,4)$, a combination of the
first two columns.)
\end{example}

In summary, elimination on an $n\times n$ matrix proceeds as follows.
\begin{enumerate}
  \item For $k=1,\dots,n$: if the $(k,k)$ entry is zero, exchange row $k$ with a lower row
    that has a nonzero entry in column $k$; if there is none, stop: $A$ is singular.
  \item For each row $i>k$, compute $\ell_{ik}=a_{ik}/a_{kk}$ (current entries) and subtract
    $\ell_{ik}$ times row $k$ from row $i$, in the augmented column as well.
  \item When all $n$ pivots are found, solve $U\vb{x}=\vb{c}$ by back substitution.
\end{enumerate}

\begin{trap}
The multipliers use the \emph{current} entries, not the original ones: in step $(3,2)$ the
multiplier is $4/2$, where $2$ is the pivot created by step $(2,1)$, not the original entry
$8$. And a zero in a pivot position is not yet a verdict of singularity: look below it first.
Only a zero with nothing nonzero underneath means failure.
\end{trap}

\subsection{Each elimination step is a matrix}

To put elimination into matrix language we need products of matrices. Multiplying a matrix
by a column gives a combination of its columns (Definition~\ref{def:u01-Ax}). The mirror image
is a row times a matrix.

\begin{proposition}[Row times matrix]\label{prop:u03-rowtimes}
If $\vb{y}\T=(y_1,\dots,y_m)$ is a row vector and $A$ is $m\times n$ with rows
$\vb{r}_1,\dots,\vb{r}_m$, then $\vb{y}\T A=y_1\vb{r}_1+\cdots+y_m\vb{r}_m$, a combination of
the rows of $A$.
\end{proposition}

\begin{proof}
Entry $j$ of $\vb{y}\T A$ is the row $\vb{y}\T$ times column $j$ of $A$, namely
$\sum_i y_ia_{ij}$. Entry $j$ of $\sum_i y_i\vb{r}_i$ is also $\sum_i y_ia_{ij}$, since
$a_{ij}$ is entry $j$ of row $\vb{r}_i$.
\end{proof}

For example, $\begin{bmatrix} 1 & 2 & 7\end{bmatrix}A
=1\,(1,2,1)+2\,(3,8,1)+7\,(0,4,1)=(7,46,10)$.
For a matrix $E$ with $m$ columns and a matrix $A$ with $m$ rows, the product $EA$ is the
matrix whose column $j$ is $E$ times column $j$ of $A$; its $(i,j)$ entry is row $i$ of $E$
times column $j$ of $A$. By Proposition~\ref{prop:u03-rowtimes}, \emph{row $i$ of $EA$ is
the combination of the rows of $A$ given by row $i$ of $E$}. (Section~\ref{sec:u04} collects
all the ways to see a matrix product.) So to subtract $3$ times row~1 from row~2 and leave
the other rows alone, the rows of $E$ must be $(1,0,0)$, $(-3,1,0)$ and $(0,0,1)$.

\begin{definition}[Elimination matrix]\label{def:u03-Eij}
For $i\neq j$, the \emph{elimination matrix} $E_{ij}$ with multiplier $\ell$ is the identity
matrix with the entry $-\ell$ placed in row $i$, column $j$. Multiplying on the left,
$E_{ij}A$ subtracts $\ell$ times row $j$ of $A$ from row $i$ and leaves the other rows
unchanged.
\end{definition}

\begin{example}[Elimination as matrix multiplication]\label{ex:u03-matrices}
For the example \eqref{eq:u03-system}, the two nontrivial steps are
\[
  E_{21}=\begin{bmatrix} 1 & 0 & 0\\ -3 & 1 & 0\\ 0 & 0 & 1\end{bmatrix},\qquad
  E_{32}=\begin{bmatrix} 1 & 0 & 0\\ 0 & 1 & 0\\ 0 & -2 & 1\end{bmatrix}.
\]
Then
\[
  E_{21}A=\begin{bmatrix} 1 & 2 & 1\\ 0 & 2 & -2\\ 0 & 4 & 1\end{bmatrix},\qquad
  E_{32}(E_{21}A)=\begin{bmatrix} 1 & 2 & 1\\ 0 & 2 & -2\\ 0 & 0 & 5\end{bmatrix}=U.
\]
Spot-check one entry: the $(2,3)$ entry of $E_{21}A$ is row~2 of $E_{21}$ times column~3 of
$A$, $(-3)(1)+(1)(1)+(0)(1)=-2$. The step $E_{31}$ is the identity, because $\ell_{31}=0$.
\end{example}

The whole of forward elimination is the single line $E_{32}(E_{31}(E_{21}A))=U$. Can we
multiply the $E$'s together first, to get one matrix that does everything? Yes, because
matrix multiplication is associative.

\begin{proposition}[Associative law]\label{prop:u03-assoc}
If the shapes allow the products, $(EF)A=E(FA)$.
\end{proposition}

\begin{proof}
Compare the $(i,j)$ entries, writing $e,f,a$ for the entries of $E,F,A$:
\[
  \bigl((EF)A\bigr)_{ij}=\sum_k\Bigl(\sum_l e_{il}f_{lk}\Bigr)a_{kj}
  =\sum_l e_{il}\Bigl(\sum_k f_{lk}a_{kj}\Bigr)=\bigl(E(FA)\bigr)_{ij}.
\]
Both are the sum of $e_{il}f_{lk}a_{kj}$ over all $k$ and $l$, added in a different order.
\end{proof}

So we may move parentheses freely, as long as the order of the factors is kept. For the
example, $E=E_{32}E_{21}$ does all of elimination at once:
\[
  E=\begin{bmatrix} 1 & 0 & 0\\ 0 & 1 & 0\\ 0 & -2 & 1\end{bmatrix}
  \begin{bmatrix} 1 & 0 & 0\\ -3 & 1 & 0\\ 0 & 0 & 1\end{bmatrix}
  =\begin{bmatrix} 1 & 0 & 0\\ -3 & 1 & 0\\ 6 & -2 & 1\end{bmatrix},\qquad EA=U,\quad E\vb{b}=\vb{c}.
\]
The entry $6$ appears because $3$ of row~1 went into row~2 and then $2$ of that row~2 went into
row~3. Section~\ref{sec:u05} shows how to avoid such entries by recording elimination the
other way around.

The order of the factors cannot be changed: the \emph{commutative} law fails for matrices.
Row exchanges show this vividly.

\begin{definition}[Permutation matrix]\label{def:u03-perm}
A \emph{permutation matrix} $P$ is the identity matrix with its rows reordered. Multiplying
on the left, $PA$ reorders the rows of $A$ in the same way.
\end{definition}

To exchange the two rows of a $2\times2$ matrix, exchange the rows of $I$:
\[
  \begin{bmatrix} 0 & 1\\ 1 & 0\end{bmatrix}\begin{bmatrix} a & b\\ c & d\end{bmatrix}
  =\begin{bmatrix} c & d\\ a & b\end{bmatrix},
  \qquad
  \begin{bmatrix} a & b\\ c & d\end{bmatrix}\begin{bmatrix} 0 & 1\\ 1 & 0\end{bmatrix}
  =\begin{bmatrix} b & a\\ d & c\end{bmatrix}.
\]
On the left, $P$ acts on rows; on the right, the same $P$ acts on columns (each column of $AP$
is a combination of the columns of $A$). So $PA\neq AP$ in general. In elimination we use
row operations, so the elimination and permutation matrices always multiply from the left.

\subsection{Undoing a step: inverse matrices}

Every elimination step can be undone. If $E_{21}$ subtracts $3$ times row~1 from row~2, the
step that brings row~2 back adds $3$ times row~1 to row~2:
\[
  E_{21}^{-1}E_{21}=
  \begin{bmatrix} 1 & 0 & 0\\ 3 & 1 & 0\\ 0 & 0 & 1\end{bmatrix}
  \begin{bmatrix} 1 & 0 & 0\\ -3 & 1 & 0\\ 0 & 0 & 1\end{bmatrix}
  =\begin{bmatrix} 1 & 0 & 0\\ 0 & 1 & 0\\ 0 & 0 & 1\end{bmatrix}=I .
\]
For example, $E_{21}$ turns $\vb{b}=(2,12,2)$ into $(2,6,2)$, and $E_{21}^{-1}$ turns
$(2,6,2)$ back into $(2,12,2)$. The matrix $I$ is the \emph{identity}: $I\vb{x}=\vb{x}$ for
every $\vb{x}$, the matrix analogue of the number $1$.

\begin{proposition}[Inverse of an elimination matrix]\label{prop:u03-Einv}
Let $E_{ij}(\ell)$ denote the elimination matrix with $-\ell$ in position $(i,j)$. Then
$E_{ij}(-\ell)\,E_{ij}(\ell)=I=E_{ij}(\ell)\,E_{ij}(-\ell)$. A permutation matrix that
exchanges two rows is its own inverse.
\end{proposition}

\begin{proof}
Apply $E_{ij}(-\ell)E_{ij}(\ell)$ to any matrix $M$. The first factor replaces row $i$ by
(row $i$) $-\ell\,$(row $j$) and leaves row $j$ untouched, since $j\neq i$. The second factor
adds $\ell\,$(row $j$) back to row $i$, restoring $M$. A matrix that leaves every $M$
unchanged is $I$ (take $M=I$). The other order is the same argument with $\ell$ replaced by
$-\ell$. Exchanging the same two rows twice restores every matrix.
\end{proof}

This is why elimination is legitimate: with $E$ the product of all the step matrices,
$A\vb{x}=\vb{b}$ and $U\vb{x}=\vb{c}$ have the same solutions. If $A\vb{x}=\vb{b}$ then
$EA\vb{x}=E\vb{b}$, that is, $U\vb{x}=\vb{c}$. Conversely, each step can be reversed, so
multiplying $U\vb{x}=\vb{c}$ on the left by the inverse steps returns $A\vb{x}=\vb{b}$. The
answer $\vb{x}=(2,1,-2)$ found by back substitution is therefore the solution of the original
system.

\begin{trap}
Row operations multiply on the \emph{left}; a matrix on the right of $A$ combines columns.
Writing $AE_{21}$ instead of $E_{21}A$ subtracts $3$ times column~2 from column~1, a
different operation entirely. Also keep the factors in order: elimination is $E_{32}E_{21}A$,
not $E_{21}E_{32}A$.
\end{trap}

\subsection{Exercises}

\begin{exercise}\label{ex:u03-two}
For the system $2x+3y=5$, $6x+15y=12$: what multiple of equation~1 should be subtracted from
equation~2? Write the system as $A\vb{x}=\vb{b}$, eliminate (what are the pivots?), and solve by
back substitution. Check your answer in both equations.
\end{exercise}
\begin{solution}
Subtract $\ell_{21}=6/2=3$ times equation~1. In matrix form,
\[
  \begin{bmatrix} 2 & 3\\ 6 & 15\end{bmatrix}\begin{bmatrix} x\\ y\end{bmatrix}
  =\begin{bmatrix} 5\\ 12\end{bmatrix}
  \;\longrightarrow\;
  \begin{bmatrix} 2 & 3\\ 0 & 6\end{bmatrix}\begin{bmatrix} x\\ y\end{bmatrix}
  =\begin{bmatrix} 5\\ -3\end{bmatrix},
\]
since $15-3\cdot3=6$ and $12-3\cdot5=-3$. The pivots are $2$ and $6$. Back substitution:
$6y=-3$ gives $y=-\tfrac12$; then $2x-\tfrac32=5$ gives $x=\tfrac{13}{4}$. Check:
$2\cdot\tfrac{13}{4}+3(-\tfrac12)=\tfrac{13}{2}-\tfrac32=5$ and
$6\cdot\tfrac{13}{4}+15(-\tfrac12)=\tfrac{39}{2}-\tfrac{15}{2}=12$.
\end{solution}

\begin{exercise}\label{ex:u03-four}
Solve by elimination, using a row exchange when needed:
\[
  x-y-z+u=0,\qquad 2x+2z=8,\qquad -y-2z=-8,\qquad 3x-3y-2z+4u=7 .
\]
\end{exercise}
\begin{solution}
The augmented matrix and the elimination steps:
\[
  \left[\begin{array}{rrrr|r} 1&-1&-1&1&0\\ 2&0&2&0&8\\ 0&-1&-2&0&-8\\ 3&-3&-2&4&7\end{array}\right]
  \to
  \left[\begin{array}{rrrr|r} 1&-1&-1&1&0\\ 0&2&4&-2&8\\ 0&-1&-2&0&-8\\ 0&0&1&1&7\end{array}\right]
  \to
  \left[\begin{array}{rrrr|r} 1&-1&-1&1&0\\ 0&2&4&-2&8\\ 0&0&0&-1&-4\\ 0&0&1&1&7\end{array}\right].
\]
The first arrow subtracts $2$ times row~1 from row~2 and $3$ times row~1 from row~4. The
second adds $\tfrac12$ of row~2 to row~3 (multiplier $\ell_{32}=-\tfrac12$). Now the $(3,3)$
position holds $0$, but row~4 has a $1$ below it, so exchange rows~3 and~4; the matrix is
then upper triangular with pivots $1,2,1,-1$. Back substitution: $-u=-4$ gives $u=4$;
$z+u=7$ gives $z=3$; $2y+4z-2u=8$ gives $2y=8-12+8=4$, $y=2$; $x-y-z+u=0$ gives $x=2+3-4=1$.
So $(x,y,z,u)=(1,2,3,4)$. Check the last equation: $3-6-6+16=7$.
\end{solution}

\begin{exercise}\label{ex:u03-pascal}
Find the lower triangular $E$ that reduces the $4\times4$ Pascal matrix to a smaller Pascal
matrix:
\[
  E\begin{bmatrix} 1&0&0&0\\ 1&1&0&0\\ 1&2&1&0\\ 1&3&3&1\end{bmatrix}
  =\begin{bmatrix} 1&0&0&0\\ 0&1&0&0\\ 0&1&1&0\\ 0&1&2&1\end{bmatrix}.
\]
Which matrix $M$, a product of several such $E$'s, reduces the Pascal matrix all the way to $I$?
\end{exercise}
\begin{solution}
Subtracting each row from the row below it uses Pascal's rule
$\binom{i}{j}-\binom{i-1}{j}=\binom{i-1}{j-1}$ and produces the smaller Pascal matrix:
\[
  E=\begin{bmatrix} 1&0&0&0\\ -1&1&0&0\\ 0&-1&1&0\\ 0&0&-1&1\end{bmatrix}.
\]
(Here row $i$ is replaced by row $i$ minus the \emph{original} row $i-1$; as a sequence of
elimination steps, do row~4 first, then row~3, then row~2.) The same idea one level down
clears the second column with $E_2$ (rows~3 and~4 minus the rows above them, leaving rows~1,
2 alone), and then $E_3$ subtracts row~3 from row~4:
\[
  E_2=\begin{bmatrix} 1&0&0&0\\ 0&1&0&0\\ 0&-1&1&0\\ 0&0&-1&1\end{bmatrix},\quad
  E_3=\begin{bmatrix} 1&0&0&0\\ 0&1&0&0\\ 0&0&1&0\\ 0&0&-1&1\end{bmatrix},\quad
  M=E_3E_2E=\begin{bmatrix} 1&0&0&0\\ -1&1&0&0\\ 1&-2&1&0\\ -1&3&-3&1\end{bmatrix}.
\]
Since $M$ times the Pascal matrix is $I$, $M$ is its inverse: Pascal's triangle with
alternating signs.
\end{solution}

\begin{exercise}\label{ex:u03-fail}
For which numbers $d$ does elimination on
$A=\begin{bmatrix} 2&1&0\\ 4&d&1\\ 0&1&2\end{bmatrix}$ meet a zero in a pivot position? For
each such $d$, say whether a row exchange rescues it.
\end{exercise}
\begin{solution}
The first pivot is $2$; subtracting $2$ times row~1 from row~2 gives row~2 $=(0,\,d-2,\,1)$.
If $d=2$ the $(2,2)$ position is zero, but row~3 $=(0,1,2)$ has a $1$ below it; exchanging
rows~2 and~3 gives pivots $2,1,1$ (the matrix becomes $\begin{bmatrix}2&1&0\\0&1&2\\0&0&1\end{bmatrix}$),
so this failure is only temporary. If $d\neq2$, the second pivot is $d-2$; the multiplier is
$\ell_{32}=1/(d-2)$, and the third pivot is $2-\frac{1}{d-2}$. It is zero when $d-2=\tfrac12$,
that is $d=\tfrac52$, and then nothing lies below it: complete failure, and $A$ is singular.
For every other $d$, elimination succeeds without exchanges.
\end{solution}
